% Generated from locale/tr/content/proof-theory/normalization/normalization-thm.tex; do not hand-edit.


\olfileid[tr]{pt}{nor}{thm}

\olsection{Normalleştirme Teoremi}

!!^a{proof}, kesme kesiti yoksa \emph{normal} biçimdedir. Bunun, ispata ne
bir permütasyon dönüştürmesi ne de indirgeme dönüştürmesi uygulanabilmesiyle
denk olduğu açıktır. !!a{proof}{}ı \emph{normalleştirmek}, hiç kesme kalmayana
dek permütasyon ve indirgeme dönüştürmeleri uygulayarak onu normal bir ispata
dönüştürmektir. Bunun her zaman mümkün olduğunu \emph{normalleştirme teoremi}
söyler.

\begin{thm}\ollabel{thm:normalization} Her $\Log{N1i}$-!!{proof}
$\delta$ normalleşir; yani kesmesiz !!a{proof}~$\delta'$ biçimine
dönüştürülebilir.
\end{thm}

\begin{proof}
  $\delta$'nın kesme derecesi ve uzunluğu üzerinde tümevarımla. Tümevarım
  varsayımı, daha düşük kesme derecesine ya da aynı kesme derecesiyle daha
kısa kesme uzunluğuna sahip her !!{proof}{}ın normalleştiğidir.

  Kesme uzunluğu~$0$ ise kesme yoktur ve $\delta$ zaten normaldir. Bu nedenle
  $\delta$'nın kesme derecesiyle kesme uzunluğunun~$>0$ olduğunu varsayın. En
  yukarıdaki, en sağdaki maksimal kesme kesitini seçin. Kesme kesitinin
  uzunluğu $>1$ ise bir permütasyon dönüştürmesi; uzunluğu~$1$ ise karşılık
  gelen indirgeme dönüştürmesini uygulayın. Böylece ya aynı kesme derecesine
  fakat daha kısa kesme uzunluğuna ya da daha düşük kesme derecesine sahip
  yeni bir !!{proof} elde edilir ve tümevarım varsayımı uygulanır.
\end{proof}

Yukarıdaki ispatta kullanılan strateji (her zaman en yukarıdaki, en sağdaki
maksimal kesmeyi ele almak), normal bir !!{proof} elde etmek için
dönüştürmelerin !!a{proof}{}a uygulanması gereken belirli bir sıra önerir.
Şaşırtıcı biçimde sıra aslında önemli değildir. Permütasyon ve indirgeme
dönüştürmelerini herhangi bir sırada uygulamak sonunda normal bir ispat verir.

Kesit ve kesme tanımları $\Log{N2i}$ için kolayca formüle edilebilir;
$\Log{N1i}$ permütasyon ve indirgeme dönüştürmeleri de doğrudan $\Log{N2i}$'ye
aktarılır. Dolayısıyla normalleştirme teoremi $\Log{N2i}$ için de geçerlidir.

\begin{thm}\ollabel{thm:normalization-N2i} Her $\Log{N2i}$-!!{proof}
$\delta$ normalleşir; yani kesmesiz !!a{proof}~$\delta'$ biçimine
dönüştürülebilir.
\end{thm}

